\chapter{Currents}

Before we discuss magnetic fields,
let's first define their sources: \textbf{currents}.

The type of current that you're probably most familiar with is the
\textbf{line current} or \textbf{filamentary current}.
It's given by
\begin{equation}
    \vec{I} = \lambda \vec{v},
\end{equation}
where \(\lambda\) is the linear current density
and \(\vec{v}\) is the velocity for the carrier particles.
\footnote{
    Some authors prefer to define line current as a purely scalar quantity
    rather than a vector quantity.
}
The unit of line current is the ampere (symbol \unit{\A}),
which is equivalent to \unit{\C \per \s}.
Note that amperes are a base unit in the SI system
and coulombs are a derived unit.

The \textbf{surface current} is given by
\begin{equation}
    \vec{K} = \sigma \vec{v},
\end{equation}
which has units of \unit{\A \per \m}.

The \textbf{volume current} or \textbf{current density} is given by
\begin{equation}
    \vec{J} = \rho \vec{v},
\end{equation}
which has units of \unit{\A \per \m\squared}.

\section{Biot-Savart}